\documentclass[11pt,a4paper]{article}

\usepackage[utf8]{inputenc}
\usepackage[T1]{fontenc}
\usepackage{amsmath,amssymb,amsfonts}
\usepackage[french]{babel}
\frenchbsetup{StandardLists=true}
\usepackage{enumitem}
\usepackage{parskip}
\usepackage[left=2cm,right=2cm,top=2.5cm,bottom=2cm]{geometry}
\usepackage{multirow,tabularx,booktabs}
\usepackage[dvipsnames]{xcolor}
\usepackage{fouriernc}
\usepackage{fancyhdr}
\usepackage{colortbl}
\usepackage{pifont}
\usepackage{chemmacros}
\chemsetup{modules=all}

%------ En-tête / pied de page ------
\pagestyle{fancy}
\renewcommand\headrulewidth{1pt}
\fancyhead[L]{Lycée Denis Diderot}
\fancyhead[R]{Classe de Terminale}
\fancyfoot[C]{page \thepage}
\fancyfoot[R]{\today}
\fancyfoot[L]{Alexis RUIZ}

%------ Commandes maison ------
\definecolor{repcolor}{RGB}{180,0,0}
\newcommand{\rep}[1]{\textcolor{repcolor}{\textbf{#1}}}

% Titre d'exercice sans tcolorbox
\newcommand{\titreexo}[2]{%
  \vspace{4mm}%
  \noindent\colorbox{BrickRed}{\color{white}\bfseries\strut\enspace #1\enspace}%
  \hfill\textbf{#2}%
  \vspace{1mm}\par\noindent\textcolor{BrickRed}{\rule{\linewidth}{1pt}}\par
  \vspace{2mm}%
}

\renewcommand{\arraystretch}{1.5}

% ============================================================
\begin{document}

\begin{center}
  \Large\rep{GRILLE DE CORRECTION}\\[2mm]
  \large Évaluation \textit{-- Prévoir une réaction d'oxydo-réduction --} Terminale\\[4mm]
  \normalsize
 
\end{center}

\vspace{2mm}\hrule\vspace{3mm}

% ============================================================
\titreexo{Exercice 1 \,--\, Cases à cocher}{/ 6 points}

% --- Question ①
\textbf{\ding{202}}\enspace Affirmations vraie / fausse
\hfill\small\textit{0,5 pt par bonne réponse \,--\, sous-total : \textbf{2,5 pts}}

\medskip
\begin{tabularx}{\linewidth}{|>{\raggedright\arraybackslash}X|>{\centering\arraybackslash}p{2.3cm}|>{\centering\arraybackslash}p{2.3cm}|>{\centering\arraybackslash}p{1.8cm}|}
\hline
\rowcolor{BrickRed!15}
\textbf{Affirmation} & \textbf{Vraie} & \textbf{Fausse} & \textbf{Pts}\\
\hline
Une réaction d'oxydo-réduction est un échange de protons
  & & \rep{\ding{52}} & 0,5\\
\hline
Une oxydation est une perte d'électrons
  & \rep{\ding{52}} & & 0,5\\
\hline
Une réduction est un gain d'électrons
  & \rep{\ding{52}} & & 0,5\\
\hline
Un oxydant perd des électrons
  & & \rep{\ding{52}} & 0,5\\
\hline
Le métal cuivre est un réducteur
  & \rep{\ding{52}} & & 0,5\\
\hline
\multicolumn{3}{|r|}{\textbf{Sous-total \ding{202}}}
  & \rep{\textbf{2,5}}\\
\hline
\end{tabularx}

\smallskip

% --- Question ②
\textbf{\ding{203}}\enspace Ions \ch{Cu\pch[2]} et métal zinc
\hfill\small\textit{0,5 pt par bonne réponse \,--\, sous-total : \textbf{2 pts}}

\medskip
\begin{tabularx}{\linewidth}{|>{\raggedright\arraybackslash}X|>{\centering\arraybackslash}p{2.3cm}|>{\centering\arraybackslash}p{2.3cm}|>{\centering\arraybackslash}p{1.8cm}|}
\hline
\rowcolor{BrickRed!15}
\textbf{Proposition} & \textbf{Oui} & \textbf{Non} & \textbf{Pts}\\
\hline
L'ion \ch{Cu\pch[2]} est l'oxydant le plus fort
  & \rep{\ding{52}} & & 0,5\\
\hline
L'ion \ch{Zn\pch[2]} est l'oxydant le plus fort
  & & \rep{\ding{52}} & 0,5\\
\hline
Le métal cuivre est le réducteur le plus fort
  & & \rep{\ding{52}} & 0,5\\
\hline
Le métal zinc est le réducteur le plus fort
  & \rep{\ding{52}} & & 0,5\\
\hline
\multicolumn{3}{|r|}{\textbf{Sous-total \ding{203}}}
  & \rep{\textbf{2,0}}\\
\hline
\end{tabularx}

\smallskip

% --- Question ③
\textbf{\ding{204}}\enspace Cocher les équations exactes \,--\, réaction \ch{Cu\pch[2]} / \ch{Fe}
\hfill\small\textit{0,5 pt par équation correcte cochée \,--\, sous-total : \textbf{1,5 pts}}

\medskip
\begin{tabularx}{\linewidth}{|>{\raggedright\arraybackslash}X|>{\centering\arraybackslash}p{3.2cm}|>{\centering\arraybackslash}p{1.8cm}|}
\hline
\rowcolor{BrickRed!15}
\textbf{Équation} & \textbf{À cocher ?} & \textbf{Pts}\\
\hline
\ch{Cu\pch[2] + 2 e\mch[] -> Cu}
  & \rep{\ding{52}\enspace Oui} & 0,5\\
\hline
\ch{Fe\pch[2] + 2 e\mch[] -> Fe}
  & \ding{55}\enspace Non & --\\
\hline
\ch{Cu\pch[2] + Fe -> Cu + Fe\pch[2]}
  & \rep{\ding{52}\enspace Oui} & 0,5\\
\hline
\ch{Cu -> Cu\pch[2] + 2 e\mch[]}
  & \ding{55}\enspace Non & --\\
\hline
\ch{Fe -> Fe\pch[2] + 2 e\mch[]}
  & \rep{\ding{52}\enspace Oui} & 0,5\\
\hline
\ch{Cu + Fe\pch[2] -> Cu\pch[2] + Fe}
  & \ding{55}\enspace Non & --\\
\hline
\multicolumn{2}{|r|}{\textbf{Sous-total \ding{204}}}
  & \rep{\textbf{1,5}}\\
\hline
\end{tabularx}

\newpage 

% ============================================================
\titreexo{Exercice 2 \,--\, QCM}{/ 7 points}

\textit{1 pt par bonne réponse. La bonne réponse est indiquée en \rep{rouge}.}

\medskip
\begin{tabularx}{\linewidth}{|>{\centering\arraybackslash}p{5mm}|>{\raggedright\arraybackslash}p{4.5cm}|>{\raggedright\arraybackslash}X|>{\centering\arraybackslash}p{1.5cm}|}
\hline
\rowcolor{BrickRed!15}
\textbf{N°} & \textbf{Bonne réponse} & \textbf{Justification / commentaire} & \textbf{Pts}\\
\hline
\ding{202}
  & \rep{Gagne des électrons}
  & L'oxydant se réduit en gagnant des électrons
  & 1\\
\hline
\ding{203}
  & \rep{Image B :} \rep{\ch{Ag} dans \ch{Cu\pch[2]}}
  & \ch{Cu\pch[2]/Cu} moins oxydant que \ch{Ag\pch[]/Ag} $\Rightarrow$ \ch{Cu\pch[2]} ne peut pas oxyder \ch{Ag} $\Rightarrow$ aucun dépôt
  & 1\\
\hline
\ding{204}
  & \rep{Un gain d'électrons}
  & Définition de la réduction
  & 1\\
\hline
\ding{205}
  & \rep{Il se réduit}
  & \ch{Ag\pch[] + e\mch[] -> Ag} : gain d'électrons = réduction
  & 1\\
\hline
\ding{206}
  & \rep{Gagnent deux \ch{e\mch[]}}
  & Demi-équation : \ch{Cu\pch[2] + 2 e\mch[] -> Cu}
  & 1\\
\hline
\ding{207}
  & \rep{\ch{Al + Cr\pch[3] -> Cr + Al\pch[3]}}
  & \ch{Cr\pch[3]} plus oxydant $\Rightarrow$ il oxyde \ch{Al} (règle du $\gamma$)
  & 1\\
\hline
\ding{208}
  & \rep{$\times 2$ la 1\textsuperscript{re}, $\times 3$ la 2\textsuperscript{e}, puis additionner}
  & PPCM(3,2)~=~6 $\Rightarrow$ $2\,\ch{Al}+3\,\ch{Zn\pch[2]} \longrightarrow 2\,\ch{Al\pch[3]}+3\,\ch{Zn}$
  & 1\\
\hline
\multicolumn{3}{|r|}{\textbf{Total Exercice 2}}
  & \rep{\textbf{7}}\\
\hline
\end{tabularx}

% ============================================================
\titreexo{Exercice 3 \,--\, Béchers}{/ 7 points}

\textit{Barème détaillé par expérience.}

\medskip
\begin{tabularx}{\linewidth}{|>{\centering\arraybackslash}p{3.2cm}|>{\raggedright\arraybackslash}X|>{\centering\arraybackslash}p{1.5cm}|}
\hline
\rowcolor{BrickRed!15}
\textbf{Expérience} & \textbf{Éléments attendus} & \textbf{Pts}\\

%--- Exp 1 ---
\hline
\multirow{2}{3.2cm}{\centering\textbf{Exp~\ding{202}}\\[2pt]\small\ch{Ag} dans \ch{Cu\pch[2]}}
  & Réponse : \rep{la réaction n'a pas lieu} & 0,5\\
\cline{2-3}
  & Justification : \ch{Cu\pch[2]/Cu} est moins oxydant que \ch{Ag\pch[]/Ag}, donc \ch{Cu\pch[2]} ne peut pas oxyder \ch{Ag} & 0,5\\


%--- Exp 2 ---
\hline
\multirow{4}{3.2cm}{\centering\textbf{Exp~\ding{203}}\\[2pt]\small\ch{Zn} dans \ch{Pt\pch[2]}}
  & Réponse : \rep{la réaction a lieu} (\ch{Pt\pch[2]} plus oxydant que \ch{Zn\pch[2]}) & 0,5\\
\cline{2-3}
  & Demi-éq. de réduction : \rep{\ch{Pt\pch[2] + 2 e\mch[] -> Pt}} & 0,5\\
\cline{2-3}
  & Demi-éq. d'oxydation : \rep{\ch{Zn -> Zn\pch[2] + 2 e\mch[]}} & 0,5\\
\cline{2-3}
  & Équation bilan (électrons s'annulent directement) : \rep{\ch{Pt\pch[2] + Zn -> Pt + Zn\pch[2]}} & 0,5\\

%--- Exp 3 ---
\hline
\multirow{4}{3.2cm}{\centering\textbf{Exp~\ding{204}}\\[2pt]\small\ch{Al} dans \ch{Sn\pch[2]}}
  & Réponse : \rep{la réaction a lieu} (\ch{Sn\pch[2]} plus oxydant que \ch{Al\pch[3]}) & 0,5\\
\cline{2-3}
  & Demi-éq. de réduction ($\times 3$) : \rep{\ch{Sn\pch[2] + 2 e\mch[] -> Sn}} & 0,5\\
\cline{2-3}
  & Demi-éq. d'oxydation ($\times 2$) : \rep{\ch{Al -> Al\pch[3] + 3 e\mch[]}} & 0,5\\
\cline{2-3}
  & Équation bilan (PPCM(2,3)~=~6) : \rep{$3\,\ch{Sn\pch[2]} + 2\,\ch{Al} \longrightarrow 3\,\ch{Sn} + 2\,\ch{Al\pch[3]}$} & 0,5\\

%--- Exp 4 ---
\hline
\multirow{4}{3.2cm}{\centering\textbf{Exp~\ding{205}}\\[2pt]\small\ch{Mg} dans \ch{Fe\pch[2]}}
  & Réponse : \rep{la réaction a lieu} (\ch{Fe\pch[2]} plus oxydant que \ch{Mg\pch[2]}) & 0,5\\
\cline{2-3}
  & Demi-éq. de réduction : \rep{\ch{Fe\pch[2] + 2 e\mch[] -> Fe}} & 0,5\\
\cline{2-3}
  & Demi-éq. d'oxydation : \rep{\ch{Mg -> Mg\pch[2] + 2 e\mch[]}} & 0,5\\
\cline{2-3}
  & Équation bilan (électrons s'annulent directement) : \rep{\ch{Fe\pch[2] + Mg -> Fe + Mg\pch[2]}} & 0,5\\
\hline

\multicolumn{2}{|r|}{\textbf{Total Exercice 3}} & \rep{\textbf{7,0}}\\
\hline
\end{tabularx}

\vspace{6mm}



\end{document}